% ============================================================================
% Курсова работа, ОКС „Магистър“, ТУ София
% Дисциплина: MCSPC05 Видео и аудио анализ в криминалистиката (семестър I)
% Програма:   Киберсигурност и превенция на киберпрестъпността (MCSPC)
%
% Тема: Възстановяване на изтрити и повредени видеофайлове при криминалистична
%       експертиза: структура на MP4 и H.264, методи и доказателствена стойност.
%
% Компилация:  latexmk -pdf main.tex        (pdfLaTeX, НЕ XeLaTeX)
% Форматиране: preamble.tex (форматът на ФКСТ, взет от Waypoint)
% Незапълнените места са маркирани с \TODO{...}:
%     grep -rn 'TODO' chapters/ main.tex
% ============================================================================
\documentclass[12pt]{article}

% ============================================================================
% preamble.tex , форматиране по изискванията на ФКСТ.
%
% Взето от waypoint-communication-networks/docs/preamble.tex (формат на ФКСТ).
% Махнати са само макросите, специфични за Waypoint.
%
% Двигател: pdfLaTeX (T2A + tempora). НЕ XeLaTeX , вж. REQUIREMENTS.md Б.0.
% ============================================================================

\usepackage{cmap}
\usepackage[T2A]{fontenc}
\usepackage[utf8]{inputenc}
\usepackage[bulgarian]{babel}

% --- Типография (А.4) ---------------------------------------------------
% tempora = Times-метрично съвместим шрифт с пълна кирилица под T2A.
\usepackage{tempora}
\usepackage[top=2.5cm, bottom=2.5cm, left=3.0cm, right=2.0cm, a4paper]{geometry}
\usepackage{setspace}
\setstretch{1.5}                     % „междуредие 1.5 lines“ , калибруем, вж. Б.2
% „без пренасяния на срички“ (А.4) , чрез ЗАБРАНИТЕЛНИ НАКАЗАНИЯ, не чрез
% \usepackage[none]{hyphenat}. Разликата е съществена: `hyphenat[none]` задава
% \hyphenchar=-1 и така изключва самия механизъм за прекъсване, от който зависи
% `breaklines` на listings , дългите редове в листингите тогава НЕ се пренасят и
% излизат извън полето (1227 препълвания дойдоха оттам). Наказанията постигат същия
% резултат за текста (TeX никога не пренася), но оставят механизма достъпен.
\usepackage[expansion=false]{microtype}
\DisableLigatures{encoding=*,family=tt*}
\hyphenpenalty=10000
\exhyphenpenalty=10000
\emergencystretch=1.5em              % цената на justified-без-преноси, вж. Б.4
\sloppy
\usepackage{indentfirst}             % отстъп и на ПЪРВИЯ абзац след заглавие, вж. Б.3
\setlength{\parindent}{1.25cm}
\setlength{\parskip}{0pt}

% --- Номерация на страниците: долу вдясно (А.4) -------------------------
\usepackage{fancyhdr}
\pagestyle{fancy}
\fancyhf{}
\renewcommand{\headrulewidth}{0pt}
\fancyfoot[R]{\thepage}

% --- Бележки под линия: 10 pt, single, justified (А.4) -------------------
\usepackage[hang,flushmargin]{footmisc}
\usepackage{ragged2e}
\renewcommand{\footnotelayout}{\fontsize{10}{12}\selectfont\setstretch{1.0}\justifying}

% --- Заглавия на раздели / параграфи / подпараграфи (А.7) ---------------
\usepackage{titlesec}
\newcommand{\sectionbreak}{\clearpage}   % „всеки раздел започва на нова страница“

\renewcommand{\thesection}{\Roman{section}}
\renewcommand{\thesubsection}{\arabic{section}.\arabic{subsection}}
\renewcommand{\thesubsubsection}{\arabic{section}.\arabic{subsection}.\arabic{subsubsection}}

% Раздел: римски, ALL CAPS, центриран, bold 14 pt
\titleformat{\section}[block]
  {\fontsize{14}{16.8}\bfseries\centering}{\thesection. }{0pt}{\MakeUppercase}

% Параграф: арабски, bold 14 pt, отстъп 1.25 cm
\titleformat{\subsection}[block]
  {\fontsize{14}{16.8}\bfseries}{\thesubsection. }{0pt}{}
\titlespacing*{\subsection}{1.25cm}{*3}{*2}

% Подпараграф: bold 13 pt, отстъп 2.0 cm
\titleformat{\subsubsection}[block]
  {\fontsize{13}{15.6}\bfseries}{\thesubsubsection. }{0pt}{}
\titlespacing*{\subsubsection}{2.0cm}{*2}{*1}

% --- Съдържание: разделите с ГЛАВНИ букви, 12 pt, bold (А.6) -------------
\usepackage{etoolbox}
\usepackage[titles]{tocloft}
% NB: кръпката ТРЯБВА да е СЛЕД tocloft , пакетът предефинира \l@section и
% иначе я изтрива (тихо: няма грешка, просто разделите излизат с малки букви).
\makeatletter
\patchcmd{\l@section}{#1}{\MakeUppercase{#1}}{}{}
\makeatother
\renewcommand{\cftsecleader}{\cftdotfill{\cftdotsep}}
\renewcommand{\cftsubsecleader}{\cftdotfill{\cftdotsep}}
\renewcommand{\cftsubsubsecleader}{\cftdotfill{\cftdotsep}}
\renewcommand{\cftsecfont}{\fontsize{12}{14}\bfseries}
\renewcommand{\cftsecpagefont}{\fontsize{12}{14}\bfseries}
\renewcommand{\cftsecaftersnum}{.\enspace}
\settowidth{\cftsecnumwidth}{VIII.\enspace}
\renewcommand{\cftsubsecaftersnum}{.}
\renewcommand{\cftsubsubsecaftersnum}{.}

% --- Надписи: таблици ОТГОРЕ, фигури ОТДОЛУ (А.8, А.9) ------------------
\usepackage{caption}
\usepackage{subcaption}
\DeclareCaptionLabelFormat{tablelabel}{Таблица #2}
\captionsetup[table]{labelformat=tablelabel, labelsep=period, position=above,
                     font={bf,small}, justification=centering, singlelinecheck=true}
\DeclareCaptionLabelFormat{figurelabel}{Фигура #2}
\captionsetup[figure]{labelformat=figurelabel, labelsep=period, position=below,
                      font={bf,small}, justification=centering, singlelinecheck=true}

% Номера на страници в библиографията без тире-диапазон (виж README).
\newcommand{\pp}{-}

% Препратки в текста: „(Табл. 3)“ / „(Фиг. 7)“ , задължителни по А.8/А.9.
\newcommand{\tabref}[1]{(Табл.~\ref{#1})}
\newcommand{\figref}[1]{(Фиг.~\ref{#1})}

% --- Математика ---------------------------------------------------------
\usepackage{amsmath,amssymb,amsfonts}
\usepackage{amsthm}

% --- Графики ------------------------------------------------------------
\usepackage{graphicx}
\usepackage{float}
\usepackage{tikz}
\usetikzlibrary{decorations.pathreplacing,arrows.meta,positioning,shapes.geometric,fit,calc}
\usepackage{adjustbox}
\usepackage{pdfpages}     % за подвързаните задание и декларация, вж. Б.5

% --- Таблици ------------------------------------------------------------
\usepackage{array}
% Ляво подравнена колона с фиксирана ширина. Тесните `p{}` колони с ДВУСТРАННО
% подравняване (наследено от основния текст) разкъсват редовете с огромни
% междудумни интервали , правилото за justified от А.4 се отнася за ОСНОВНИЯ текст,
% не за клетки от 3 cm. Всички таблици с текстови клетки ползват L.
\newcolumntype{L}[1]{>{\RaggedRight\arraybackslash}p{#1}}
\usepackage{booktabs}
\usepackage[table]{xcolor}% \cellcolor за топлинната карта
\usepackage{longtable}    % „продължение“ при пренасяне на таблица (А.8)
\usepackage{multirow}
\usepackage{tabularx}

% --- Помощни ------------------------------------------------------------
\usepackage{hyperref}
\usepackage{xurl}
\usepackage{enumitem}
\usepackage{textcomp}
\usepackage{listings}
\urlstyle{same}
% КРИТИЧНО: `breaklines` на listings вмъква ДИСКРЕЦИОННИ прекъсвания, а глобалните
% \hyphenpenalty/\exhyphenpenalty=10000 (правилото „без пренасяния“, А.4) ги правят
% безкрайно скъпи , тогава дългите редове в листингите НЕ се пренасят и излизат извън
% полето. Проверено изолирано: с наказанията 1 препълване, без тях 0. Наказанията се
% нулират САМО вътре в листинг, така че текстът си остава без пренасяния.
\makeatletter
\lst@AddToHook{PreInit}{\hyphenpenalty=50\exhyphenpenalty=50}
\makeatother
\renewcommand{\lstlistingname}{Листинг}
\renewcommand{\lstlistlistingname}{Списък на листингите}

% Italic САМО за ключови понятия (А.4) , семантичен макрос, за да е ревизуемо.
\newcommand{\term}[1]{\emph{#1}}
% Инлайн код/идентификатори.
\DeclareRobustCommand{\code}[1]{\texttt{#1}}
% Видим маркер за незапълнено място. Търси се с: grep -rn 'TODO' chapters/
\newcommand{\TODO}[1]{\textcolor{red}{\textbf{[TODO: #1]}}}

\tikzset{
    block/.style={draw, rounded corners, fill=blue!5, align=center, font=\small,
        minimum width=3.2cm, minimum height=1.05cm, inner sep=6pt},
    accent/.style={draw, rounded corners, fill=green!8, align=center, font=\small,
        minimum width=3.0cm, minimum height=0.95cm, inner sep=6pt},
    danger/.style={draw, rounded corners, fill=red!6, align=center, font=\small,
        minimum width=3.0cm, minimum height=0.95cm, inner sep=6pt},
    note/.style={draw, dashed, rounded corners, fill=gray!8, align=left, font=\small,
        inner sep=6pt},
    line/.style={-Latex, thick}
}

% --- Единично разстояние в нетекстови среди -----------------------------
\AtBeginEnvironment{tikzpicture}{\setstretch{1.1}}
\AtBeginEnvironment{lstlisting}{\setstretch{1.1}}
\AtBeginEnvironment{verbatim}{\setstretch{1.1}}
\AtBeginEnvironment{tabular}{\setstretch{1.1}}
\AtBeginEnvironment{longtable}{\setstretch{1.1}}

\BeforeBeginEnvironment{tikzpicture}{\begin{adjustbox}{max width=\textwidth}}
\AfterEndEnvironment{tikzpicture}{\end{adjustbox}}

% --- Оцветяване на C++ код ----------------------------------------------
\definecolor{cppblue}{RGB}{0,0,180}
\definecolor{cppgreen}{RGB}{0,128,0}
\definecolor{cppred}{RGB}{180,0,0}
\definecolor{cppgray}{RGB}{128,128,128}

\lstset{
    language=C++,
    basicstyle=\footnotesize\ttfamily,
    % Пренасяне на дълги редове. `breakatwhitespace=false` е СЪЩЕСТВЕНО: коментарите в
    % изходния текст стигат 110 знака без удобно място за прекъсване, и при `true` редът
    % изобщо не се пренася, а излиза извън полето. `linewidth` се задава изрично, защото
    % рамката и левият отстъп иначе не се приспадат от наличната ширина.
    breaklines=true, breakatwhitespace=false, breakautoindent=false, breakindent=0pt,
    linewidth=\dimexpr\linewidth-2.5em\relax,
    columns=flexible, keepspaces=true,
    postbreak=\mbox{\textcolor{cppgray}{\ensuremath{\hookrightarrow}}\space},
    xleftmargin=1.5em, xrightmargin=0.5em,
    keywordstyle=\color{cppblue}\bfseries,
    stringstyle=\color{cppred},
    commentstyle=\color{cppgreen},
    identifierstyle=\color{black},
    showstringspaces=false,
    numbers=left, numberstyle=\tiny\color{cppgray}, numbersep=8pt,
    aboveskip=1.5em, belowskip=1.5em,
    frame=single, framerule=0.5pt, rulecolor=\color{cppgray},
    tabsize=4, captionpos=b,
    morekeywords={concept, requires, constexpr, consteval, constinit, noexcept,
                  nullptr, override, final, co_await, co_return, co_yield,
                  expected, unexpected, span, string_view, unique_ptr, shared_ptr},
    emph={maui, view, element, handler, layout, page, application, bindable_object,
          view_ref, weak_ref, observable, grid, stack_layout, content_page},
    emphstyle=\color{cppblue},
    inputencoding=utf8,
    % Не-ASCII знаци, срещащи се в коментарите на изходния текст. Без тези редове
    % \lstinputlisting отказва с „Invalid UTF-8 byte sequence“ , listings не разбира
    % многобайтов вход, който не е в literate таблицата.
    literate={→}{{$\rightarrow$}}1 {×}{{$\times$}}1 {…}{{...}}1 {§}{{\S}}1
{а}{{\cyra}}1 {б}{{\cyrb}}1 {в}{{\cyrv}}1 {г}{{\cyrg}}1 {д}{{\cyrd}}1
             {е}{{\cyre}}1 {ж}{{\cyrzh}}1 {з}{{\cyrz}}1 {и}{{\cyri}}1 {й}{{\cyrishrt}}1
             {к}{{\cyrk}}1 {л}{{\cyrl}}1 {м}{{\cyrm}}1 {н}{{\cyrn}}1 {о}{{\cyro}}1
             {п}{{\cyrp}}1 {р}{{\cyrr}}1 {с}{{\cyrs}}1 {т}{{\cyrt}}1 {у}{{\cyru}}1
             {ф}{{\cyrf}}1 {х}{{\cyrh}}1 {ц}{{\cyrc}}1 {ч}{{\cyrch}}1 {ш}{{\cyrsh}}1
             {щ}{{\cyrshch}}1 {ъ}{{\cyrhrdsn}}1 {ь}{{\cyrsftsn}}1 {ю}{{\cyryu}}1 {я}{{\cyrya}}1
             {А}{{\CYRA}}1 {Б}{{\CYRB}}1 {В}{{\CYRV}}1 {Г}{{\CYRG}}1 {Д}{{\CYRD}}1
             {Е}{{\CYRE}}1 {Ж}{{\CYRZH}}1 {З}{{\CYRZ}}1 {И}{{\CYRI}}1 {Й}{{\CYRISHRT}}1
             {К}{{\CYRK}}1 {Л}{{\CYRL}}1 {М}{{\CYRM}}1 {Н}{{\CYRN}}1 {О}{{\CYRO}}1
             {П}{{\CYRP}}1 {Р}{{\CYRR}}1 {С}{{\CYRS}}1 {Т}{{\CYRT}}1 {У}{{\CYRU}}1
             {Ф}{{\CYRF}}1 {Х}{{\CYRH}}1 {Ц}{{\CYRC}}1 {Ч}{{\CYRCH}}1 {Ш}{{\CYRSH}}1
             {Щ}{{\CYRSHCH}}1 {Ъ}{{\CYRHRDSN}}1 {Ь}{{\CYRSFTSN}}1 {Ю}{{\CYRYU}}1 {Я}{{\CYRYA}}1,
}


\begin{document}

% ============================================================================
% Заглавна страница
% ============================================================================
\begin{titlepage}
\thispagestyle{empty}
\begin{center}
    {\fontsize{18}{21.6}\selectfont\bfseries ТЕХНИЧЕСКИ УНИВЕРСИТЕТ, СОФИЯ\par}
    \vspace{0.5cm}
    {\fontsize{16}{19.2}\selectfont\bfseries ФАКУЛТЕТ „КОМПЮТЪРНИ СИСТЕМИ И ТЕХНОЛОГИИ“\par}

    \vspace{0.3cm}
    \begin{figure}[H]
        \centering
        \includegraphics[width=0.30\linewidth]{images/tu-logo.png}
    \end{figure}
    \vspace{0.2cm}

    {\fontsize{20}{23}\selectfont\bfseries ВЪЗСТАНОВЯВАНЕ НА ИЗТРИТИ И ПОВРЕДЕНИ ВИДЕОФАЙЛОВЕ ПРИ КРИМИНАЛИСТИЧНА ЕКСПЕРТИЗА\par}

    \vspace{0.6cm}

    {\fontsize{14}{16.8}\selectfont\bfseries КУРСОВА РАБОТА\par}
    \vspace{0.3cm}
    {\fontsize{14}{16.8}\selectfont\bfseries по дисциплина „Видео и аудио анализ в криминалистиката“\par}

    \vfill

    \begin{flushleft}
    \begin{tabular}{@{}l@{\hspace{0.6em}}l@{}}
    Студент: & Алекс Иванов Цветанов\\
    Група: & 263\\
    Фак. №: & \TODO{фак. номер}\\
    \end{tabular}

    \vspace{0.7cm}

    % Дисциплината се води от трима преподаватели. Кой е поставил заданието,
    % още не е записано. Да се впише само той.
    {\fontsize{14}{16.8}\selectfont\bfseries Преподавател:\par}
    \TODO{име на преподавателя}\par
    \end{flushleft}

    \vfill

    {\fontsize{12}{14.4}\selectfont София, 2026\par}
\end{center}
\end{titlepage}

% ============================================================================
% Съдържание
% ============================================================================
\newpage
\pagestyle{empty}
\setcounter{tocdepth}{2}
\tableofcontents
\newpage

\pagestyle{fancy}
\pagenumbering{arabic}

% ============================================================================
% УВОД, без номер
% ============================================================================
\section*{УВОД}
\addcontentsline{toc}{section}{УВОД}

Видеозаписът е едно от най-често срещаните веществени доказателства в съвременното
разследване. Източниците са много: камери за видеонаблюдение, видеорегистратори в
автомобили, нагръдни камери, мобилни телефони. Общото между тях е, че записът
съществува като файл върху цифров носител. Затова преди да започне анализът на
съдържанието (лица, регистрационни номера, последователност на събития), експертът
трябва да разполага с файл, който може да бъде прочетен и възпроизведен.

На практика това често не е така. Файлът може да е изтрит от потребителя, носителят
може да е форматиран, записът може да е прекъснат от спиране на захранването, а
части от файла може да са презаписани с нови данни. В тези случаи стандартните
програми за възпроизвеждане отказват да отворят файла или показват само част от него.
Възстановяването на такива файлове е отделна задача на цифровата криминалистика, която
стои преди същинския видеоанализ.

\term{Целта на курсовата работа} е да опише теоретичната основа на възстановяването на
изтрити и повредени видеофайлове и да оцени доказателствената стойност на резултата.
За постигането ѝ са поставени следните задачи:

\begin{enumerate}[leftmargin=2cm]
  \item да се опише структурата на видеофайл във формат MP4 и на видеопотока H.264,
        доколкото тя определя какво може да бъде възстановено;
  \item да се систематизират причините за загуба и повреда на видеофайлове;
  \item да се сравнят основните методи за възстановяване: по метаданни на файловата
        система, чрез изрязване по сигнатури (\term{file carving}), за фрагментирани
        файлове, на ниво кадри и чрез ремонт с еталонен файл;
  \item да се опишат изискванията към процедурата, така че възстановеният запис да
        запази доказателствената си стойност.
\end{enumerate}

Работата е теоретична. Тя не съдържа програмен код и собствени измервания. Всички
технически твърдения се позовават на стандартите и на публикуваните източници,
посочени в списъка с литература.

% ============================================================================
% I. Видеофайлът като обект на експертиза
% ============================================================================
\section{Видеофайлът като обект на експертиза}

Един видеофайл съдържа две различни нива на структура. Външното ниво е
\term{контейнерът}: той описва какви потоци има във файла (видео, звук, субтитри),
къде във файла се намира всяка порция данни и кога трябва да бъде възпроизведена.
Вътрешното ниво е \term{кодекът}: той определя как самото изображение е компресирано
в поредица от байтове. За възстановяването двете нива имат различно значение. Загубата
на контейнера прави файла нечетим за обикновен плейър, но компресираните кадри може
да са още налице. Загубата на част от кодираните данни поврежда самите кадри.

В работата е разгледана най-разпространената комбинация: контейнер MP4 и видеокодек
H.264 (наричан още AVC). Изборът не е случаен. MP4 и сродният му формат MOV се
използват масово от мобилни телефони и камери, а H.264 е кодекът, с който работи
голяма част от записващите устройства.

\subsection{Контейнерът MP4 и базовият медиен формат на ISO}

Форматът MP4 \cite{iso14496_14} е частен случай на \term{базовия медиен файлов формат
на ISO} (ISO Base Media File Format, ISOBMFF), дефиниран в ISO/IEC 14496-12
\cite{iso14496_12}. Файлът е изграден от последователност от \term{кутии} (boxes, в
по-старите текстове \term{atoms}). Всяка кутия започва със заглавие от два елемента:
размер от 4 байта и тип от 4 знака, например \code{ftyp}, \code{moov} или \code{mdat}.
Ако размерът е 1, след типа следва 64-битов размер. Ако размерът е 0, кутията
продължава до края на файла. Кутиите могат да се влагат една в друга, така че файлът
образува дърво.

Три кутии от най-горно ниво са важни за възстановяването \figref{fig:mp4}:

\begin{itemize}[leftmargin=2cm]
  \item \code{ftyp} (file type) посочва към кой вариант на формата принадлежи файлът
        и се поставя възможно най-рано в него. Когато е първата кутия, типът
        \code{ftyp} стои на байтове 4 до 7 от началото на файла и служи като
        \term{сигнатура} при търсене на изтрити файлове.
  \item \code{mdat} (media data) съдържа самите компресирани данни на видеото и
        звука. Това обикновено е най-голямата част от файла.
  \item \code{moov} (movie) съдържа метаданните: описание на всяка писта
        (\code{trak}) и таблиците на пробите (\code{stbl}).
\end{itemize}

\begin{figure}[H]
\centering
\begin{tikzpicture}[node distance=0.35cm]
  \node[block, minimum width=3.0cm] (ftyp) {\code{ftyp}\\ тип на файла};
  \node[accent, right=of ftyp, minimum width=4.6cm] (moov) {\code{moov}\\ метаданни};
  \node[danger, right=of moov, minimum width=4.6cm] (mdat) {\code{mdat}\\ компресирани данни};

  \node[note, below=0.9cm of moov.south east, anchor=north east, text width=6.6cm] (stbl)
    {\code{trak} $\rightarrow$ \code{mdia} $\rightarrow$ \code{minf} $\rightarrow$ \code{stbl}:\\
     \code{stsd} (параметри на кодека, SPS и PPS)\\
     \code{stsz} (размер на всяка проба)\\
     \code{stco} / \code{co64} (отместване на всяко парче)\\
     \code{stss} (кои проби са ключови кадри)};
  \node[note, right=1.4cm of stbl, text width=4.0cm] (nal)
    {Пробите на видеото:\\
     {}[дължина][NAL единица]\\
     {}[дължина][NAL единица]\\
     \ldots};

  \draw[line] (moov.south) to (moov.south |- stbl.north);
  \draw[line] (mdat.south) to (mdat.south |- nal.north);
  \draw[line, dashed] (stbl.east) to node[above, font=\scriptsize]{адреси} (nal.west);
\end{tikzpicture}
\caption{Опростена структура на файл MP4 (по \cite{iso14496_12, iso14496_15})}
\label{fig:mp4}
\end{figure}

Ключовото наблюдение е, че кутията \code{mdat} сама по себе си не е самоописваща се.
Тя е непрекъснат поток от байтове. Къде започва и свършва всеки кадър, личи само от
таблиците в \code{moov}: \code{stsz} дава размера на всяка проба, \code{stco} или
\code{co64} дава отместването на всяко парче (chunk) от началото на файла, а
\code{stsc} свързва пробите с парчетата \cite{iso14496_12}. Параметрите, нужни за
декодиране на видеото, също са в \code{moov}, в описанието на пробите \code{stsd}.
Ако \code{moov} липсва, плейърът не знае нито къде са кадрите, нито как да ги декодира.

Стандартът предвижда и \term{фрагментиран} вариант, при който метаданните са
разпределени в множество кутии \code{moof}, всяка от които описва следващия я участък
от данни \cite{iso14496_12}. Така файлът остава четим до последния завършен фрагмент.
Значението на тази разлика за повредени записи е разгледано в раздел~II.

\subsection{Видеопотокът H.264}

Стандартът H.264 \cite{itu_h264} дели кодирания поток на \term{единици на слоя за
мрежова абстракция} (Network Abstraction Layer units, NAL). Всяка NAL единица започва
с 1 байт заглавие. В него 5 бита задават типа на единицата, 2 бита (\code{nal\_ref\_idc})
показват дали тя служи за опора при декодиране на други кадри, а 1 бит трябва винаги
да е нула. Таблица~\ref{tab:nal} изброява типовете, които имат значение за
възстановяването.

\begin{table}[H]
\caption{Типове NAL единици в H.264, важни за възстановяването (по \cite{itu_h264})}
\label{tab:nal}
\centering
\begin{tabular}{@{}cL{4.2cm}L{7.0cm}@{}}
\toprule
\textbf{Тип} & \textbf{Наименование} & \textbf{Значение за възстановяването} \\
\midrule
1 & Част от кадър, който не е IDR & Декодира се само ако е наличен кадърът, от който зависи \\
5 & Част от IDR кадър & Декодира се самостоятелно, след него зависимостите се нулират \\
6 & SEI & Допълнителна информация, не е нужна за декодиране \\
7 & SPS & Параметри на поредицата: размер на кадъра, профил \\
8 & PPS & Параметри на картината, сочи към SPS \\
\bottomrule
\end{tabular}
\end{table}

Компресията в H.264 използва зависимости между кадрите \cite{wiegand2003}. Кадрите от
тип I се кодират само с информация от самите тях. Кадрите от тип P и B се кодират като
разлика спрямо други, вече декодирани кадри. Специален вид I кадър е \term{IDR кадърът}
(Instantaneous Decoding Refresh): след него нито един следващ кадър не се позовава на
кадрите преди него. Следствието за експертизата е пряко. Ако е загубен IDR кадър, всички
зависими от него кадри до следващия IDR кадър не могат да бъдат декодирани коректно.
Обратно, всеки оцелял IDR кадър е отправна точка, от която декодирането може да
започне наново.

Стандартът описва и два начина за разграничаване на NAL единиците в поток. В
\term{приложение B} на H.264 всяка единица се предшества от \term{стартов код}, байтовете
\code{00 00 01} \cite{itu_h264}. Във файл MP4 се използва друг запис: всяка единица се
предшества от полето ѝ за дължина, а стартови кодове няма \cite{iso14496_15}. Размерът
на това поле, както и SPS и PPS, се пазят в кутията \code{avcC} вътре в \code{stsd}, т.е.
отново в \code{moov}. Затова видеоданни от MP4 без \code{moov} не могат да се намерят
с просто търсене на стартови кодове. Методите в раздел~III трябва да разпознаят
полетата за дължина и да възстановят параметрите на кодека по друг път.

% ============================================================================
% II. Причини за загуба и повреда
% ============================================================================
\section{Причини за загуба и повреда на видеофайлове}

Методът за възстановяване се избира според това какво точно се е случило с файла.
Различните причини оставят различни следи върху носителя. По-долу те са групирани в
пет случая.

\subsection{Изтриване на файла}

При изтриване файлът обикновено не се унищожава. Файловата система само отбелязва
като изтрит записа, който сочи къде се намира файлът. Данните остават на носителя, но
операционната система вече смята мястото за свободно и може да го презапише по всяко
време \cite{nist80086}. На магнитен твърд диск данните остават физически на носителя, докато
мястото не бъде използвано за нов запис. Доколко могат да се възстановят и метаданните
(име, размер, списък на клъстерите), зависи от конкретната файлова система и от това колко
време е работил носителят след изтриването.

При твърдотелните дискове (SSD) това предположение не е сигурно. Bell и Boddington
показват експериментално, че вграденият софтуер на SSD може сам да изтрие
освободените блокове чрез т.нар. \term{събиране на отпадъци} (garbage collection), без
команда от компютъра \cite{bell2010}. Тогава изтритият видеозапис може да е унищожен
още преди носителят да бъде иззет.

\subsection{Форматиране на носителя}

Бързото форматиране създава нови, празни структури на файловата система. На магнитен
твърд диск то не изтрива самите данни: в опита на Bell и Boddington бързо форматиран
твърд диск запазва всички файлове \cite{bell2010}. Връзката между файловете и техните клъстери обаче
се губи. Тогава остава единствено търсенето на файлове по тяхното съдържание, без
помощта на файловата система (раздел~3.2). Ограничението от раздел~2.1 важи и тук. В същия опит
бързо форматиран SSD изчиства почти всички данни сам, за около три минути
\cite{bell2010}.

\subsection{Прекъснат запис}

Този случай е специфичен за видеото. Устройството може да пише файла MP4
последователно: първо \code{ftyp}, после растящата кутия \code{mdat}, а кутията
\code{moov} едва при нормално спиране на записа, когато таблиците на пробите са вече
известни. Ако при такава подредба записът бъде прекъснат (изтощена батерия, отнето
захранване, извадена карта, повреда на устройството), \code{moov} не се записва. Резултатът е файл с пълни видеоданни, но без метаданни.
Обикновените плейъри не могат да го отворят, въпреки че почти цялото съдържание е
налице. За възстановяване на такива прекъснати записи е създаден инструментът \code{untrunc}
\cite{untrunc},
разгледан в раздел~3.5. Фрагментираният вариант на MP4 (раздел~1.1) намалява загубата,
защото всеки завършен фрагмент носи собствените си метаданни.

\subsection{Фрагментиране на файла върху носителя}

Файловата система не винаги разполага файла в последователни клъстери. Когато свободното
място е разпокъсано, един файл може да бъде разделен на две или повече несъседни части
\cite{garfinkel2007}. Докато метаданните на файловата система са налице, това не е проблем.
След изтриване или форматиране обаче фрагментираният файл не може да бъде възстановен
чрез просто копиране на байтовете от началото до края. Garfinkel изследва над 300 употребявани
носителя и намира, че около 6\,\% от файловете на тях са фрагментирани. Разпределението
е силно неравномерно: на около половината носители няма нито един фрагментиран файл, а
на 30 носителя делът надхвърля 10\,\% \cite{garfinkel2007}. Самият автор определя
фрагментирането като сравнително рядко, но смята възстановяването на фрагментирани файлове
за необходима възможност на всеки инструмент \cite{garfinkel2007}.

\subsection{Частично презаписване}

Когато част от клъстерите на изтрит файл е използвана за нови данни, съответната част
от записа е загубена безвъзвратно. Никой метод не може да възстанови презаписаните
байтове. Възможно е обаче да се спаси останалото. Na и съавтори показват, че при
повреден видеофайл с презаписани участъци по-голямата част от съдържанието в
непрезаписаните участъци може да бъде възстановена \cite{na2014}. Тук значение има
структурата на H.264 от раздел~1.2: оцелелите кадри могат да се декодират, ако е оцелял
и кадърът, от който зависят.

Таблица~\ref{tab:causes} обобщава какво остава на носителя във всеки от петте случая.

\begin{table}[H]
\caption{Какво остава на носителя при различните причини за загуба}
\label{tab:causes}
\centering
\begin{tabular}{@{}L{3.6cm}L{3.4cm}L{3.0cm}L{3.6cm}@{}}
\toprule
\textbf{Причина} & \textbf{Метаданни на файловата система} & \textbf{Кутия \code{moov}} & \textbf{Видеоданни} \\
\midrule
Изтриване & Частично & Налична & Налични, ако не са презаписани \\
Форматиране & Загубени & Налична & Налични, ако не са презаписани \\
Прекъснат запис & Налични & Липсва & Налични \\
Фрагментиране & Загубени (след изтриване) & Може да е в друг фрагмент & Разпокъсани \\
Презаписване & Без значение & Може да е загубена & Частично загубени \\
\bottomrule
\end{tabular}
\end{table}

% ============================================================================
% III. Методи за възстановяване
% ============================================================================
\section{Методи за възстановяване}

Методите са подредени от най-простия към най-сложния. Всеки следващ метод разчита на
по-малко информация от носителя и затова изисква повече знание за самия формат. На
практика експертът започва от първия метод и преминава към следващите само ако
предишният не е дал резултат.

\subsection{Възстановяване по метаданни на файловата система}

При изтриване записът, който сочи къде е файлът, обикновено само се отбелязва като
изтрит \cite{nist80086}. Ако този запис и данните са оцелели, файлът може да бъде
възстановен по информацията на самата файлова система. Това е
най-надеждният случай: възстановява се точно последователността от клъстери, която е
съставлявала файла, включително при фрагментиране. Ограничението е, че методът зависи
изцяло от оцелелите структури на файловата система. След форматиране или след
продължителна работа на носителя тези структури обикновено са загубени или
презаписани.

\subsection{Изрязване по сигнатури (file carving)}

\term{Изрязването на файлове} (file carving) е възстановяване на файлове само по тяхното
съдържание, без помощта на файловата система \cite{pal2009}. Програмата обхожда целия
носител и търси известни \term{сигнатури}: поредици от байтове, с които започват
(заглавие) или завършват (окончание) файловете от даден формат. Между намереното
заглавие и следващото окончание (или до определен максимален размер) байтовете се
изрязват като кандидат за файл. Инструментът Scalpel е пример за такъв подход, при който
сигнатурите се задават в конфигурационен файл, а носителят се обхожда ефективно
\cite{richard2005scalpel}.

За MP4 подходящата сигнатура е типът \code{ftyp} на байтове 4 до 7 от началото на файла,
когато \code{ftyp} е първата кутия (раздел~1.1). Окончание в същия смисъл MP4 няма. Структурата на кутиите обаче дава
по-добра възможност: от заглавието на всяка кутия се чете нейният размер, така че
програмата може да „прескача“ от кутия на кутия и да определи точния край на файла.
Това е частен случай на \term{валидиране}, разгледано в следващата точка.

Основното ограничение на изрязването по сигнатури е предположението, че файлът е записан
непрекъснато. При фрагментиран файл изрязаният участък съдържа чужди данни, а
действителното продължение на файла е другаде \cite{garfinkel2007}.

\subsection{Възстановяване на фрагментирани файлове}

Garfinkel предлага два подхода за този проблем \cite{garfinkel2007}. Първият е
\term{бързото валидиране на обекти} (fast object validation): всеки кандидат за файл се
проверява дали е формално валиден за своя формат, например дали се декодира без грешки.
Валидирането отхвърля погрешно изрязаните кандидати и позволява да се изпробват различни
комбинации от участъци. Вторият подход е \term{изрязването с откриване на празнина при
два фрагмента} (bifragment gap carving). Ако файлът е разделен на две части с чужди данни
между тях, програмата изпробва различни размери и положения на празнината, докато
валидирането не потвърди, че резултатът е цял файл.

Pal и Memon проследяват развитието на изрязването като цяло \cite{pal2009}. Общият
принцип на методите за фрагментирани файлове е един и същ: колкото по-строго може да се провери, че
даден участък е продължение на файла, толкова по-надеждно е сглобяването. За видео в
MP4 проверката може да използва две нива: дали кутиите имат съгласувани размери и дали
кодираните данни в \code{mdat} се декодират като валиден поток H.264.

\subsection{Възстановяване на ниво кадри}

Когато контейнерът е повреден или липсва, а част от данните е презаписана, търсенето на
цял файл губи смисъл. Na и съавтори предлагат друг подход: възстановяване на отделни
кадри по спецификацията на кодека \cite{na2014}. Методът разпознава кодираните кадри
директно в данните и ги декодира, като използва знанието за структурата на потока. Така
се спасява съдържанието от непрезаписаните участъци дори когато файлът като цяло не
може да бъде отворен. Както посочват авторите, методът възстановява по-голямата част от
непрезаписаното съдържание, но не и презаписаните участъци \cite{na2014}.

Този подход се опира пряко на структурата от раздел~1.2. За да се декодират кадрите, са
нужни SPS и PPS. Ако те са загубени заедно с \code{moov}, параметрите трябва да се
възстановят от друг източник, например от запис на същото устройство. Освен това
декодирането може да започне само от IDR кадър. Кадрите между началото на оцелелия
участък и първия IDR кадър в него остават недекодируеми.

\subsection{Ремонт на прекъснат запис с еталонен файл}

Прекъснатият запис от раздел~2.3 е специален случай: видеоданните са налице и
непрекъснати, липсва само \code{moov}. Инструментът \code{untrunc} решава точно тази задача
\cite{untrunc}. Той изисква \term{еталонен файл}: друг, неповреден запис от същата камера
или приложение \cite{untrunc}. Защо е нужен еталон, следва от раздел~I. На повредения
файл липсват параметрите на кодека и таблиците на пробите, а те са в \code{moov}.
Параметрите на кодека могат да се вземат от еталон, записан със същите настройки.
Размерите и отместванията на пробите трябва да се изчислят наново от самите данни в
\code{mdat}. Така по принцип се получава нова кутия \code{moov}, която прави файла
възпроизводим. Документацията на инструмента не описва вътрешната му работа и
предупреждава, че за успеха е нужен „и малко късмет“ \cite{untrunc}.

\subsection{Сравнение на методите}

Таблица~\ref{tab:methods} обобщава петте метода. Изборът между тях следва от
таблица~\ref{tab:causes}: за всяка причина за загуба има метод, който използва точно
онова, което е оцеляло.

\begin{table}[H]
\caption{Сравнение на методите за възстановяване}
\label{tab:methods}
\centering
\begin{tabular}{@{}L{3.3cm}L{3.6cm}L{3.4cm}L{3.6cm}@{}}
\toprule
\textbf{Метод} & \textbf{Какво изисква} & \textbf{Какво възстановява} & \textbf{Основно ограничение} \\
\midrule
По метаданни на файловата система & Оцелели структури на файловата система & Целия файл, вкл. фрагментиран & Не работи след форматиране \\
Изрязване по сигнатури & Само данните на носителя & Цели непрекъснати файлове & Греши при фрагментиране \\
Валидиране и сглобяване & Строга проверка на формата & Фрагментирани файлове & Изчислително скъпо \\
На ниво кадри & Параметри на кодека & Отделни кадри & Не възстановява презаписаното \\
Ремонт с еталон & Еталонен запис от същото устройство & Възпроизводим файл & Зависи от еталона \\
\bottomrule
\end{tabular}
\end{table}

% ============================================================================
% IV. Процедура и доказателствена стойност
% ============================================================================
\section{Процедура и доказателствена стойност}

Възстановеният видеозапис има стойност за разследването само ако съдът може да се
довери, че той отразява съдържанието на иззетия носител и не е създаден или изменен
при самата експертиза. Затова техническите методи от раздел~III трябва да се прилагат
в рамките на процедура, която може да бъде проверена. Стандартът ISO/IEC 27037 дава
насоки за идентифициране, събиране, придобиване и съхраняване на цифрови доказателства
\cite{iso27037}. Ръководството на NIST SP 800-86 описва процеса на цифровата
криминалистика в четири етапа: събиране, изследване, анализ и докладване
\cite{nist80086}.

\subsection{Работа върху копие}

Първото правило е, че оригиналният носител не се променя. Възстановяването се извършва
върху побитово копие (образ) на носителя, а не върху самия носител
\cite{nist80086}. Това е особено важно при възстановяване, защото всяка
операция за запис върху оригинала, дори случайна, може да презапише точно онези
освободени клъстери, които съдържат изтрития видеозапис.

\subsection{Контролни суми и верига на съхранение}

Целостта на образа се доказва с криптографска контролна сума (хеш), изчислена при
придобиването. Повторното изчисляване по всяко време трябва да даде същия резултат. Това
позволява на независим експерт да провери, че работата е извършена върху същите данни
\cite{nist80086}. Всяко действие с носителя и с образа се документира: кой, кога и с
каква цел го е извършил. Тази документация образува \term{веригата на съхранение}
(chain of custody) \cite{nist80086}.

Контролната сума се изчислява и за всеки възстановен файл. Така в доклада може да се
посочи точно кой файл е обект на последващия видеоанализ и къде в образа на носителя
се намират данните, от които е получен.

\subsection{Възстановено и реконструирано съдържание}

Методите от раздел~III се различават по това колко от резултата идва от носителя и
колко е добавено при експертизата. Тази разлика трябва да е видима в доклада.

\begin{itemize}[leftmargin=2cm]
  \item При възстановяване по метаданни на файловата система и при изрязване на
        непрекъснат файл резултатът е копие на байтове, които физически са на носителя.
        Твърдението е проверимо: всеки байт може да бъде намерен в образа.
  \item При сглобяване на фрагментиран файл байтовете също са от носителя, но
        тяхната подредба е решение на експерта или на програмата. Трябва да се посочи кои
        участъци са съединени и как е проверено, че те принадлежат на един файл.
  \item При ремонт с еталонен файл новата кутия \code{moov} не идва от носителя. Тя е
        построена от програмата въз основа на друг файл. Видеоданните са оригинални, но
        тяхното разделяне на кадри и времето на възпроизвеждане са изчислени. Това трябва
        да се отбележи изрично, заедно с произхода на еталонния файл.
  \item При възстановяване на ниво кадри се възстановяват отделни кадри, а не файлът.
        Липсващите кадри и участъци трябва да се посочат, за да не се създава впечатление
        за непрекъснат запис.
\end{itemize}

Същото важи и за времевите данни. Времето на всеки кадър във файл MP4 се изчислява от
таблиците в \code{moov} \cite{iso14496_12}. Ако \code{moov} е построена наново,
изчисленото време на кадрите е резултат от реконструкцията, а не оригинален запис.
Изводите за продължителност на събитие или за поредност на действия трябва да отчитат
това.

\subsection{Ограничения}

Методите имат и общи граници, които експертът трябва да заяви. Презаписаните данни са
загубени и никой метод не ги връща \cite{na2014}. При твърдотелни дискове освободените
блокове може да са изчистени от самия носител, така че изрязването да не намери нищо
\cite{bell2010}. Съвременните мобилни телефони освен това често криптират данните, и
без ключа възстановените байтове не могат да бъдат прочетени. Липсата на кадри не може да се тълкува
като липса на събития в записания интервал. Накрая, успешното декодиране доказва само,
че данните са валиден видеопоток. То не доказва, че записът не е бил редактиран преди
изтриването. Проверката за автентичност на съдържанието е отделна задача на
видеоанализа и излиза извън обхвата на тази работа.

% ============================================================================
% ЗАКЛЮЧЕНИЕ, без номер
% ============================================================================
\section*{ЗАКЛЮЧЕНИЕ}
\addcontentsline{toc}{section}{ЗАКЛЮЧЕНИЕ}

В курсовата работа е разгледано възстановяването на изтрити и повредени видеофайлове
като подготвителен етап на видеоанализа в криминалистиката. Изводите са следните.

Първо, възможностите за възстановяване се определят от структурата на формата. Във
файл MP4 видеоданните в \code{mdat} не носят сами информация за границите на кадрите. Тази
информация, както и параметрите на кодека, са в кутията \code{moov}. Затова загубата на
\code{moov} прави файла нечетим, дори когато всички кадри са налице.

Второ, всяка причина за загуба оставя различни следи на носителя. Изтриването запазва
данните, форматирането губи връзката им с файловете, прекъснатият запис губи
метаданните, фрагментирането разпокъсва файла, а презаписването унищожава част от
данните безвъзвратно.

Трето, на тези случаи съответстват различни методи: възстановяване по файловата
система, изрязване по сигнатури, сглобяване на фрагменти с валидиране, възстановяване
на ниво кадри и ремонт с еталонен файл. Колкото по-малко е оцеляло от структурата на
файла, толкова повече методът разчита на знание за формата и толкова повече от
резултата е изчислено, а не прочетено от носителя.

Четвърто, доказателствената стойност на резултата зависи от процедурата. Работата върху
копие, контролните суми, веригата на съхранение и ясното разграничаване между
възстановено и реконструирано съдържание са това, което позволява възстановеният запис
да бъде използван в разследването.

Работата е ограничена до формата MP4 и кодека H.264. Естествено продължение е да се
разгледат други контейнери и кодеци, както и проверката за автентичност на
възстановения запис.


% ============================================================================
% Литература
% ============================================================================
\newpage
\phantomsection
\addcontentsline{toc}{section}{ЛИТЕРАТУРА}
\bibliographystyle{ieeetr}
\bibliography{references}

\end{document}
